% Unit 051: English translation of chapter4.tex, lines 260--560.
% Source: Methods of Algebra, Volume 2, commit
% 9a5803ff77dd3257484cb177f851a73770a59dd3 (CC BY 4.0).
% stable-unit-id: o014.aljabr2.chapter4.triangulated-basics
% Coverage: Section 4.2 on the basic properties of triangulated categories;
% ends before Section 4.3.

% segment-id: o014.li2.chapter4.triangulated-basics.h001
\section{Basic Properties}\label{sec:triangular-basic}
\hypertarget{o014.aljabr2.chapter4.triangulated-basics}{ }

% segment-id: o014.li2.chapter4.triangulated-basics.p001
The preceding section gave only the definitions and initial properties of
triangulated categories and cohomological functors. This section develops
their structure further.

% segment-id: o014.li2.chapter4.triangulated-basics.pr001
\begin{proposition}\label{prop:dt-3-2}
	Consider the following commutative diagram in a pretriangulated category
	$\mathcal{D}$:
% segment-id: o014.li2.chapter4.triangulated-basics.g001
	\[\begin{tikzcd}
		X \arrow[r] \arrow[d, "\alpha"'] & Y \arrow[d, "\beta"] \arrow[r] & Z \arrow[d, "\gamma"] \arrow[r] & TX \arrow[d, "T\alpha"] \\
		X' \arrow[r] & Y' \arrow[r] & Z' \arrow[r] & TX'
	\end{tikzcd}\]
	If each row is a distinguished triangle and any two of $\alpha$, $\beta$,
	and $\gamma$ are isomorphisms, then the remaining morphism is also an
	isomorphism.
\end{proposition}

% segment-id: o014.li2.chapter4.triangulated-basics.pf001
\begin{proof}
	After suitable rotations (TR3), we may assume without loss of generality
	that $\alpha$ and $\gamma$ are isomorphisms. To show that $\beta$ is an
	isomorphism, it suffices to prove, for every object $S$, that
	$\Hom(S, Y) \xrightarrow{\beta_*} \Hom(S, Y')$ is an isomorphism.
	Proposition~\ref{prop:cohomological-Hom-functor} gives the following
	commutative diagram with exact rows:
% segment-id: o014.li2.chapter4.triangulated-basics.g002
	\[\begin{tikzcd}[column sep=small]
		\Hom(S, T^{-1} Z) \arrow[d, "\sim" sloped] \arrow[r] & \Hom(S, X) \arrow[r] \arrow[d, "\sim" sloped] & \Hom(S,Y) \arrow[d, "{\beta_*}"] \arrow[r] & \Hom(S,Z) \arrow[d, "\sim" sloped] \arrow[r] & \Hom(S, TX) \arrow[d, "\sim" sloped] \\
		\Hom(S, T^{-1} Z') \arrow[r] & \Hom(S, X') \arrow[r] & \Hom(S, Y') \arrow[r] & \Hom(S, Z') \arrow[r] & \Hom(S, TX') .
	\end{tikzcd}\]
	Now apply Proposition~\ref{prop:5-lemma}.
\end{proof}

% segment-id: o014.li2.chapter4.triangulated-basics.p002
In the argument above, $\Hom(S, \cdot)$ may also be replaced by
$\Hom(\cdot, S)$.

% segment-id: o014.li2.chapter4.triangulated-basics.r001
\begin{remark}\label{rem:special-triangles}
	For the moment, call a triangle $X \to Y \to Z \xrightarrow{+1}$ a
	\emph{special triangle} (or a \emph{cospecial triangle}) if, for every
	$S \in \Obj(\mathcal{D})$, the sequence
% segment-id: o014.li2.chapter4.triangulated-basics.q001
	\[ \cdots \to \Hom(S, X) \to \Hom(S, Y) \to \Hom(S, Z) \to \Hom(S, TX) \to \cdots \]
	(or the corresponding version with $\Hom(\cdot, S)$) is exact. If
	$X_i \to Y_i \to Z_i \xrightarrow{+1}$ is a family of distinguished
	triangles ($i \in I$), then
	$\prod_i X_i \to \prod_i Y_i \to \prod_i Z_i \xrightarrow{+1}$ (or
	$\coprod_i X_i \to \coprod_i Y_i \to \coprod_i Z_i \xrightarrow{+1}$)
	is a special triangle (or a cospecial triangle), provided the products
	(or coproducts) in question exist. This follows simply from the universal
	properties of products and coproducts together with
	\cite[Lemma~6.8.3]{Li1}. Notice that because
	$T: \mathcal{D} \to \mathcal{D}$ is an equivalence, it necessarily
	preserves products and coproducts. Proposition~\ref{prop:dt-3-2} extends
	to the case in which both rows are special (or cospecial) triangles.
\end{remark}

% segment-id: o014.li2.chapter4.triangulated-basics.co001
\begin{corollary}\label{prop:Cone-uniqueness}
	For every morphism $f: X \to Y$ in a pretriangulated category
	$\mathcal{D}$, the distinguished triangle
	$X \xrightarrow{f} Y \to Z \xrightarrow{+1}$ in axiom (TR2) is unique
	up to isomorphism.
\end{corollary}

% segment-id: o014.li2.chapter4.triangulated-basics.pf002
\begin{proof}
	Given distinguished triangles
	$X \xrightarrow{f} Y \to Z \xrightarrow{+1}$ and
	$X \xrightarrow{f} Y \to Z' \xrightarrow{+1}$, axiom (TR4) gives a
	morphism $\gamma$ that makes the diagram
% segment-id: o014.li2.chapter4.triangulated-basics.g003
	\[\begin{tikzcd}
		X \arrow[r, "f"] \arrow[d, "\identity_X"'] & Y \arrow[r] \arrow[d, "\identity_Y"] & Z \arrow[r] \arrow[d, "\gamma"] & TX \arrow[d, "\identity_{TX}"] \\
		X \arrow[r, "f"'] & Y \arrow[r] & Z' \arrow[r] & T X
	\end{tikzcd}\]
	commute. Proposition~\ref{prop:dt-3-2} then implies that $\gamma$ is an
	isomorphism.
\end{proof}

% segment-id: o014.li2.chapter4.triangulated-basics.p003
It must be emphasized that the isomorphism in
Corollary~\ref{prop:Cone-uniqueness} is \emph{not canonical}.

% segment-id: o014.li2.chapter4.triangulated-basics.le001
\begin{lemma}\label{prop:dt-sum}
	Let $\mathcal{D}$ be a pretriangulated category and let
	$X_i \xrightarrow{f_i} Y_i \xrightarrow{g_i} Z_i \xrightarrow{h_i} TX_i$
	be a family of distinguished triangles. Suppose that
	$\prod_i X_i$, $\prod_i Y_i$, and $\prod_i Z_i$ all exist. Then
% segment-id: o014.li2.chapter4.triangulated-basics.q002
	\[ \prod_{i \in I} X_i \xrightarrow{\prod_i f_i } \prod_{i \in I} Y_i \xrightarrow{\prod_i g_i} \prod_{i \in I} Z_i \xrightarrow{\prod_i h_i} \underbracket{\prod_{i \in I} TX_i}_{\simeq T(\prod_{i \in I} X_i)} \]
	is also a distinguished triangle. The corresponding assertion holds with
	$\prod_{i \in I}$ replaced by $\coprod_{i \in I}$.
\end{lemma}

% segment-id: o014.li2.chapter4.triangulated-basics.pf003
\begin{proof}
	By duality, we need only treat the case of $\prod_{i \in I}$. From now on
	we always identify $\prod_i TX_i$ with $T(\prod_i X_i)$. By (TR2), choose
	a distinguished triangle
	$\prod_i X_i \xrightarrow{\prod_i f_i} \prod_i Y_i \to Q
	\xrightarrow{h} \prod_i TX_i$. For every $i \in I$, apply (TR4) to
	obtain a commutative diagram
% segment-id: o014.li2.chapter4.triangulated-basics.g004
	\[\begin{tikzcd}
		\prod_{j \in I} X_j \arrow[r, "{\prod_j f_j}"] \arrow[d] & \prod_j Y_j \arrow[r] \arrow[d] & Q \arrow[r, "h"] \arrow[d, "{\exists \gamma_i}"] & \prod_{j \in I} TX_j \arrow[d] \\
		X_i \arrow[r, "{f_i}"'] & Y_i \arrow[r, "{g_i}"'] & Z_i \arrow[r, "{h_i}"'] & TX_i
	\end{tikzcd}\]
	in which all vertical morphisms other than $\gamma_i$ are the canonical
	projections. The universal property of the product then gives a
	commutative diagram
% segment-id: o014.li2.chapter4.triangulated-basics.g005
	\[\begin{tikzcd}
		\prod_{i \in I} X_i \arrow[r, "{\prod_i f_i}"] \arrow[d, "\identity"'] & \prod_{i \in I} Y_i \arrow[r] \arrow[d, "\identity"] & Q \arrow[r, "h"] \arrow[d, "{\prod_i \gamma_i}"] & \prod_{i \in I} TX_i \arrow[d, "\identity"] \\
		\prod_{i \in I} X_i \arrow[r, "{\prod_i f_i}"'] & \prod_{i \in I} Y_i \arrow[r, "{\prod_i g_i}"'] & \prod_{i \in I} Z_i \arrow[r, "{\prod_i h_i}"'] & \prod_{i \in I} TX_i .
	\end{tikzcd}\]
	The second row is a special triangle in the sense of
	Remark~\ref{rem:special-triangles}. The extended version of
	Proposition~\ref{prop:dt-3-2} now implies that
	$\prod_{i \in I} \gamma_i$ is an isomorphism.
\end{proof}

% segment-id: o014.li2.chapter4.triangulated-basics.co002
\begin{corollary}\label{prop:dt-sum-0}
	In every pretriangulated category, a triangle of the form
	$X \xrightarrow{\iota_1} X \oplus Y \xrightarrow{p_2} Y
	\xrightarrow{0} TX$ is always distinguished; here $\iota_i$ and $p_j$
	are defined as in \S\ref{sec:Ker-Coker}.
\end{corollary}

% segment-id: o014.li2.chapter4.triangulated-basics.pf004
\begin{proof}
	Apply Lemma~\ref{prop:dt-sum} to the distinguished triangles
	$X \xrightarrow{\identity_X} X \to 0 \to TX$ and
	$0 \to Y \xrightarrow{\identity_Y} Y \to 0$.
\end{proof}

% segment-id: o014.li2.chapter4.triangulated-basics.p004
Conversely, if one morphism in a distinguished triangle is $0$, then the
triangle comes from a direct sum. This fact can be verified by rotating the
triangle into the following form.

% segment-id: o014.li2.chapter4.triangulated-basics.pr002
\begin{proposition}\label{prop:dt-sum-1}
	In every pretriangulated category, a distinguished triangle of the form
	$X \to M \to Y \xrightarrow{0} TX$ must be isomorphic to the triangle
	$X \to X \oplus Y \to Y \xrightarrow{0} TX$ from
	Corollary~\ref{prop:dt-sum-0}.
\end{proposition}

% segment-id: o014.li2.chapter4.triangulated-basics.pf005
\begin{proof}
	A suitably rotated version of (TR4) shows that there is a morphism
	$\alpha: X \oplus Y \to M$ such that the following diagram is a morphism
	of distinguished triangles:
% segment-id: o014.li2.chapter4.triangulated-basics.g006
	\[\begin{tikzcd}
		X \arrow[r, "\iota_1"] \arrow[d, "\identity_X"'] & X \oplus Y \arrow[r, "p_2"] \arrow[d, "\alpha"] & Y \arrow[r, "0"] \arrow[d, "\identity_Y"] & TX \arrow[d, "\identity_{TX}"] \\
		X \arrow[r] & M  \arrow[r] & Y \arrow[r, "0"']  & TX \\
	\end{tikzcd}\]
	Proposition~\ref{prop:dt-3-2} implies that $\alpha$ is an isomorphism.
\end{proof}

% segment-id: o014.li2.chapter4.triangulated-basics.p005
The next result shows that additivity in the definition of a triangulated
functor is actually redundant.

% segment-id: o014.li2.chapter4.triangulated-basics.co003
\begin{corollary}\label{prop:triangulated-automatic-additivity}
	Let $\mathcal{D}$ and $\mathcal{D}'$ be pretriangulated categories, with
	both translation functors denoted by $T$. Let
	$F: (\mathcal{D}, T) \to (\mathcal{D}', T)$ be a functor between
	categories with translation. If $F$ maps distinguished triangles to
	distinguished triangles, then $F$ is additive and hence is automatically
	a triangulated functor.
\end{corollary}

% segment-id: o014.li2.chapter4.triangulated-basics.pf006
\begin{proof}
	The distinguished triangle
	$0 \xrightarrow{\identity} 0 \xrightarrow{\identity} 0
	\xrightarrow{\identity} 0$ in $\mathcal{D}$ gives the distinguished
	triangle
	$F(0) \xrightarrow{\identity} F(0) \xrightarrow{\identity} F(0)
	\xrightarrow{\identity} F(0)$ in $\mathcal{D}'$.
	Lemma~\ref{prop:dt-composite-null} then gives
	$\identity_{F(0)} = \identity_{F(0)} \circ \identity_{F(0)} = 0$, so
	$F(0) \simeq 0$. This also implies that $F$ maps zero morphisms to zero
	morphisms.

	Next consider the distinguished triangle in $\mathcal{D}$ supplied by
	Corollary~\ref{prop:dt-sum-0}, namely
	$X \xrightarrow{\iota_1} X \oplus Y \xrightarrow{p_2} Y
	\xrightarrow{0} TX$. The corresponding triangle
% segment-id: o014.li2.chapter4.triangulated-basics.q003
	\[ FX \xrightarrow{F\iota_1} F(X \oplus Y) \xrightarrow{Fp_2} FY \xrightarrow{0} TFX \]
	is distinguished in $\mathcal{D}'$. Take $M := F(X \oplus Y)$ in
	Proposition~\ref{prop:dt-sum-1}. The morphism $\alpha$ in the proof of
	that proposition can clearly be taken to be the canonical morphism
	$FX \oplus FY \to F(X \oplus Y)$ determined by the universal property of
	the direct sum. Thus $\alpha$ is an isomorphism. By
	Corollary~\ref{prop:automatic-additivity} (iii), $F$ is additive.
\end{proof}

% segment-id: o014.li2.chapter4.triangulated-basics.p006
The next two results show that, in an equivalence of pretriangulated
categories (Definition~\ref{def:triangulated-cat}), it is enough to assume
that the functor in one direction is triangulated. This is another
application of Proposition~\ref{prop:cohomological-Hom-functor}.

% segment-id: o014.li2.chapter4.triangulated-basics.pr003
\begin{proposition}\label{prop:adjoint-triangulated}
	Let $\mathcal{D}$ and $\mathcal{D}'$ be pretriangulated categories, and
	consider functors
% segment-id: o014.li2.chapter4.triangulated-basics.g007
	$\begin{tikzcd}
		\mathcal{D} \arrow[shift left, r, "F"] & \mathcal{D}' \arrow[shift left, l, "G"]
	\end{tikzcd}$
	such that $(F, G)$ is an adjunction. Then $F$ is a triangulated functor if
	and only if $G$ is a triangulated functor. In this case the unit
	$\eta: \identity_{\mathcal{D}} \to GF$ and counit
	$\varepsilon: FG \to \identity_{\mathcal{D}'}$ of the adjunction are
	compatible with the translation functors; see
	Definition~\ref{def:category-with-translation}.
\end{proposition}

% segment-id: o014.li2.chapter4.triangulated-basics.pf007
\begin{proof}
	By duality, it suffices to consider the case in which $F$ is a
	triangulated functor. Corollary~\ref{prop:automatic-additivity} shows that
	$G$ is additive. Denote both translation functors on $\mathcal{D}$ and
	$\mathcal{D}'$ by $T$. Transporting the structure gives a new adjunction
% segment-id: o014.li2.chapter4.triangulated-basics.q004
	\[ \hat{F} := T F T^{-1}, \quad \hat{G} := T G T^{-1}, \quad \hat{\eta} := T\eta T^{-1}, \quad \hat{\varepsilon} := T \varepsilon T^{-1}. \]
	Since $\hat{F} \simeq F$, uniqueness of right adjoints
	\cite[Proposition~2.6.10]{Li1} gives a corresponding isomorphism
	$\hat{G} \simeq G$ and commutative diagrams
% segment-id: o014.li2.chapter4.triangulated-basics.g008
	\[\begin{tikzcd}
		\hat{F} \hat{G} \arrow[r, "\hat{\varepsilon}"] \arrow[d, "\sim" sloped] & \identity_{\mathcal{D}'} \\
		FG \arrow[ru, "{\varepsilon}"'] &
	\end{tikzcd} \quad \begin{tikzcd}
		\hat{G} \hat{F} \arrow[d, "\sim" sloped] & \identity_{\mathcal{D}} \arrow[l, "\hat{\eta}"'] \arrow[ld, "\eta"] \\
		GF &
	\end{tikzcd}\]
	In particular, $GT \simeq TG$ and
	$FGT \simeq FTG \simeq TFG$. After composing the commutative diagram for
	$\hat{\varepsilon}$ and $\varepsilon$ on the right with $T$, it can be
	rewritten as
% segment-id: o014.li2.chapter4.triangulated-basics.g009
	\begin{equation}\label{eqn:adjoint-triangulated-aux0}\begin{tikzcd}
		TFG \arrow[r, "{T\varepsilon}"] \arrow[d, "\sim" sloped] & T \\
		FGT \arrow[ru, "\varepsilon T"'] &
	\end{tikzcd}\end{equation}
	This equation and its dual version for $\eta$ show that $\eta$ and
	$\varepsilon$ are compatible with the translation functors.

	We now show that $G$ preserves distinguished triangles. Consider a
	distinguished triangle $X \xrightarrow{f} Y \xrightarrow{g} Z
	\xrightarrow{+1}$ in $\mathcal{D}'$. Choose a distinguished triangle
	$GX \xrightarrow{Gf} GY \xrightarrow{h} A \xrightarrow{+1}$ in
	$\mathcal{D}$, apply $F$ to it, and then use (TR3) to obtain the
	commutative diagram
% segment-id: o014.li2.chapter4.triangulated-basics.g010
	\begin{equation*}\begin{tikzcd}
		FG X \arrow[r, "FGf"] \arrow[d, "\varepsilon_X"'] & FG Y \arrow[d, "\varepsilon_Y"] \arrow[r, "Fh"] & FA \arrow[r] \arrow[d] & TFGX \arrow[d, "T\varepsilon_X"] \\
		X \arrow[r, "f"'] & Y \arrow[r, "g"'] & Z \arrow[r] & TX .
	\end{tikzcd}\end{equation*}

	Apply the adjunction to this diagram, and adjust the rightmost square
	using \eqref{eqn:adjoint-triangulated-aux0} and $GT \simeq TG$. One
	readily obtains a commutative diagram
% segment-id: o014.li2.chapter4.triangulated-basics.g011
	\[\begin{tikzcd}
		GX \arrow[r, "Gf"] \arrow[d, "\identity_X"'] & GY \arrow[d, "\identity_Y"] \arrow[r, "h"] & A \arrow[r] \arrow[d] & TGX \arrow[d, "\identity_{TGX}"] \\
		GX \arrow[r, "Gf"'] & GY \arrow[r, "Gg"'] & GZ \arrow[r] & TGX
	\end{tikzcd}\]
	whose first row is a distinguished triangle. For every
	$B \in \Obj(\mathcal{D})$, apply $\Hom_{\mathcal{D}}(B, \cdot)$ to obtain
	a commutative diagram
% segment-id: o014.li2.chapter4.triangulated-basics.g012
	\[\begin{tikzcd}[column sep=small, every cell/.append style={font=\small}]
		\Hom(B, GX) \arrow[r, "(Gf)_*" inner sep=0.8em] \arrow[d, "\simeq"'] & \Hom(B, GY) \arrow[r, "h_*" inner sep=0.8em] \arrow[d, "\simeq"] & \Hom(B, A) \arrow[d] \arrow[r] & \Hom(B, TGX) \arrow[r] \arrow[d, "\simeq"] & \Hom(B, TGY) \arrow[d, "\simeq"] \\
		\Hom(B, GX) \arrow[r, "(Gf)_*"' inner sep=0.8em] & \Hom(B, GY) \arrow[r, "(Gg)_*"' inner sep=0.8em] & \Hom(B, GZ) \arrow[r] & \Hom(B, TGX) \arrow[r] & \Hom(B, TGY)
	\end{tikzcd}\]
	Proposition~\ref{prop:cohomological-Hom-functor} shows that the first row
	is exact. On the other hand, by $TG \simeq GT$ and the adjunction, the
	second row is isomorphic to the sequence induced by the distinguished
	triangle $X \xrightarrow{f} Y \xrightarrow{g} Z \xrightarrow{+1}$,
	namely
% segment-id: o014.li2.chapter4.triangulated-basics.g013
	\[\begin{tikzcd}[column sep=small, every cell/.append style = {font = \small}]
		\Hom(FB, X) \arrow[r, "f_*"] & \Hom(FB, Y) \arrow[r, "g_*"] & \Hom(FB, Z) \arrow[r] & \Hom(FB, TX) \arrow[r] & \Hom(FB, TY) ,
	\end{tikzcd}\]
	and hence is also exact. Since $B$ is arbitrary,
	Proposition~\ref{prop:5-lemma} implies that $A \to GZ$ is an isomorphism.
	Thus $GX \xrightarrow{Gf} GY \xrightarrow{Gg} GZ \to TGX$ is indeed a
	distinguished triangle in $\mathcal{D}$.
\end{proof}

% segment-id: o014.li2.chapter4.triangulated-basics.co004
\begin{corollary}\label{prop:triangulated-equivalence}
	Let $\mathcal{D}$ and $\mathcal{D}'$ be pretriangulated categories and
	let the functors
% segment-id: o014.li2.chapter4.triangulated-basics.g014
	$\begin{tikzcd}
		\mathcal{D} \arrow[shift left, r, "F"] & \mathcal{D}' \arrow[shift left, l, "G"]
	\end{tikzcd}$
	be quasi-inverse to one another. Then:
% segment-id: o014.li2.chapter4.triangulated-basics.l001
	\begin{compactenum}[(i)]
% segment-id: o014.li2.chapter4.triangulated-basics.i001
		\item $F$ is a triangulated functor if and only if $G$ is a
		triangulated functor;
% segment-id: o014.li2.chapter4.triangulated-basics.i002
		\item if $F$ is a triangulated functor, then $F$ and $G$ give an
		equivalence of the pretriangulated categories.
	\end{compactenum}
\end{corollary}

% segment-id: o014.li2.chapter4.triangulated-basics.pf008
\begin{proof}
	Use \cite[Theorem~2.6.12]{Li1} to realize $(F, G)$ and $(G, F)$ as
	adjunctions, and then apply Proposition~\ref{prop:adjoint-triangulated}.
\end{proof}

% segment-id: o014.li2.chapter4.triangulated-basics.dp001
\begin{definition-proposition}[Pretriangulated subcategories]\label{def:triangulated-subcat}
	\index{pretriangulated category!pretriangulated subcategory}
	\index{triangulated category!triangulated subcategory}
	Let $\mathcal{D}'$ be a full additive subcategory of a pretriangulated
	category $\mathcal{D}$ that is closed under the translation functor $T$.
	Denote the inclusion functor by
	$\iota: \mathcal{D}' \to \mathcal{D}$.
% segment-id: o014.li2.chapter4.triangulated-basics.l002
	\begin{enumerate}[(i)]
% segment-id: o014.li2.chapter4.triangulated-basics.i003
		\item If a pretriangulated structure on $\mathcal{D}'$ making $\iota$
		a triangulated functor exists, it is unique. In this case, a triangle
		$X \to Y \to Z \xrightarrow{+1}$ in $\mathcal{D}'$ is distinguished
		if and only if it is distinguished in $\mathcal{D}$.
% segment-id: o014.li2.chapter4.triangulated-basics.i004
		\item This pretriangulated structure exists if and only if the
		following condition holds: whenever
		$X \to Y \to Z \xrightarrow{+1}$ is a distinguished triangle in
		$\mathcal{D}$ and any two of its terms belong to
		$\Obj(\mathcal{D}')$, the remaining term is isomorphic to an object of
		$\mathcal{D}'$.
% segment-id: o014.li2.chapter4.triangulated-basics.i005
		\item If, in addition, $\mathcal{D}$ is assumed to be triangulated,
		then $\mathcal{D}'$ is also triangulated.
	\end{enumerate}

	Such a pretriangulated (or triangulated) category $\mathcal{D}'$ is called
	a \emph{pretriangulated subcategory} (or \emph{triangulated subcategory})
	of $\mathcal{D}$.
\end{definition-proposition}

% segment-id: o014.li2.chapter4.triangulated-basics.pf009
\begin{proof}
	First consider (i). Suppose that $\iota$ is triangulated. Every
	distinguished triangle in $\mathcal{D}'$ is of course distinguished in
	$\mathcal{D}$. Conversely, suppose that a distinguished triangle
	$X \xrightarrow{f} Y \to Z \xrightarrow{+1}$ in $\mathcal{D}$ satisfies
	$X, Y, Z \in \Obj(\mathcal{D}')$. Choose any distinguished triangle
	$X \xrightarrow{f} Y \to Z' \xrightarrow{+1}$ in $\mathcal{D}'$.
	Applying Corollary~\ref{prop:Cone-uniqueness} in $\mathcal{D}$ shows that
	the two triangles are isomorphic. Since $\mathcal{D}'$ is full, (TR0)
	implies that $X \xrightarrow{f} Y \to Z \xrightarrow{+1}$ is also
	distinguished in $\mathcal{D}'$. This is precisely the pretriangulated
	structure asserted in (i).

	If this pretriangulated structure exists, then
	Corollary~\ref{prop:Cone-uniqueness} shows that the condition in (ii) is
	necessary, since after rotation we may always assume that
	$X, Y \in \Obj(\mathcal{D}')$. Conversely, assume the condition in (ii)
	and define the distinguished triangles in $\mathcal{D}'$ as stated in
	(i). Axioms (TR0), (TR1), (TR3), and (TR4) for $\mathcal{D}'$ are
	inherited directly from $\mathcal{D}$. For (TR2), given a morphism
	$f: X \to Y$ in $\mathcal{D}'$, choose a distinguished triangle
	$X \xrightarrow{f} Y \to Z \xrightarrow{+1}$ in $\mathcal{D}$. The
	condition says that $Z$ is isomorphic to an object of $\mathcal{D}'$; we
	may therefore assume that $Z \in \Obj(\mathcal{D}')$ and thus obtain a
	distinguished triangle in $\mathcal{D}'$. This proves that $\mathcal{D}'$
	is pretriangulated and establishes sufficiency.

	For (iii), the point is to verify (TR5) for $\mathcal{D}'$. Indeed, the
	triangle $Z' \to Y' \to X' \xrightarrow{+1}$ obtained by applying (TR5)
	in $\mathcal{D}$ is automatically distinguished in $\mathcal{D}'$ by (i).
\end{proof}

% segment-id: o014.li2.chapter4.triangulated-basics.cv001
\begin{convention}\label{con:saturated-intersection}
	\index{subcategory!saturated (replete)}
	Fix a category $\mathcal{C}$. A subcategory $\mathcal{C}'$ is called a
	\emph{saturated subcategory} if it is closed under isomorphisms, that is,
% segment-id: o014.li2.chapter4.triangulated-basics.q005
	\[ \forall X \in \Obj(\mathcal{C}'), \; \forall Y \in \Obj(\mathcal{C}), \quad X \simeq Y \implies Y \in \Obj(\mathcal{C}') . \]
	For every nonempty family of full subcategories
	$(\mathcal{C}_i)_{i \in I}$, their intersection
	$\bigcap_{i \in I} \mathcal{C}_i$ is by definition the full subcategory
	of $\mathcal{C}$ satisfying
	$\Obj\left( \bigcap_{i \in I} \mathcal{C}_i \right)
	= \bigcap_{i \in I} \Obj(\mathcal{C}_i)$.
\end{convention}

% segment-id: o014.li2.chapter4.triangulated-basics.pr004
\begin{proposition}\label{prop:triangulated-subcat-intersection}
	Let $\mathcal{D}$ be a pretriangulated (or triangulated) category, and let
	$(\mathcal{D}_i)_{i \in I}$ be a family of saturated pretriangulated (or
	saturated triangulated) subcategories, with $I \neq \emptyset$. Then
	$\bigcap_{i \in I} \mathcal{D}_i$ is also a saturated pretriangulated (or
	saturated triangulated) subcategory.
\end{proposition}

% segment-id: o014.li2.chapter4.triangulated-basics.pf010
\begin{proof}
	An intersection of saturated subcategories plainly preserves the
	condition in Definition--Proposition~\ref{def:triangulated-subcat} (ii).
\end{proof}

% segment-id: o014.li2.chapter4.triangulated-basics.p007
A common technique is to cut out a saturated pretriangulated (or
triangulated) subcategory using a cohomological functor.

% segment-id: o014.li2.chapter4.triangulated-basics.pr005
\begin{proposition}\label{prop:cut-subcat-by-cohomologies}
	Let $\mathcal{D}$ be a pretriangulated (or triangulated) category,
	$H: \mathcal{D} \to \mathcal{A}$ a cohomological functor, and
	$\mathcal{T}$ a weak Serre subcategory of $\mathcal{A}$
	(Definition~\ref{def:Serre-subcat}). Define a full subcategory
	$\mathcal{D}_{H, \mathcal{T}}$ of $\mathcal{D}$ by
% segment-id: o014.li2.chapter4.triangulated-basics.q006
	\[ X \in \Obj(\mathcal{D}_{H, \mathcal{T}}) \iff \forall n \in \Z, \; H(T^n X) \in \Obj(\mathcal{T}). \]
	Then $\mathcal{D}_{H, \mathcal{T}}$ is a saturated pretriangulated (or
	saturated triangulated) subcategory of $\mathcal{D}$.
\end{proposition}

% segment-id: o014.li2.chapter4.triangulated-basics.pf011
\begin{proof}
	It suffices to verify the condition in
	Definition--Proposition~\ref{def:triangulated-subcat} (ii). Recall that a
	weak Serre subcategory is saturated. Hence saturation and $T$-invariance
	of $\mathcal{D}_{H, \mathcal{T}}$ are clear. Now consider a distinguished
	triangle $X \to Y \to Z \xrightarrow{+1}$ in $\mathcal{D}$, two of whose
	terms lie in $\mathcal{D}_{H, \mathcal{T}}$. After rotation, we may assume
	without loss of generality that
	$X, Z \in \Obj(\mathcal{D}_{H, \mathcal{T}})$. For every $n \in \Z$,
	there is an exact sequence
% segment-id: o014.li2.chapter4.triangulated-basics.q007
	\[ H(T^{n-1} Z) \to H(T^n X) \to H(T^n Y) \to H(T^n Z) \to H(T^{n+1} X). \]
	The definition of a weak Serre subcategory immediately gives
	$H(T^n Y) \in \Obj(\mathcal{T})$.
\end{proof}

% segment-id: o014.li2.chapter4.triangulated-basics.ex001
\begin{example}\label{eg:subcat-N-H}
	\index[sym1]{NH@$\mathcal{N}_H$}
	Let $\mathcal{D}$ be a pretriangulated (or triangulated) category and let
	$H: \mathcal{D} \to \mathcal{A}$ be a cohomological functor. Define
	$\mathcal{N}_H$ to be the full subcategory whose objects satisfy
	$H(T^n X) = 0$ for every $n \in \Z$. This is the case of
	Proposition~\ref{prop:cut-subcat-by-cohomologies} in which $\mathcal{T}$
	is the subcategory consisting of zero objects. Thus $\mathcal{N}_H$ is a
	pretriangulated (or triangulated) subcategory.
\end{example}

% segment-id: o014.li2.chapter4.triangulated-basics.p008
The following technical result will be used later in
\S~\sourcecrossref{sec:triangulated-localization}{4.3}; it depends on the
octahedral axiom (TR5).

% segment-id: o014.li2.chapter4.triangulated-basics.le002
\begin{lemma}[J.-L.\ Verdier]\label{prop:triangulated-9-diag}
	Every commutative diagram in a triangulated category $\mathcal{D}$
% segment-id: o014.li2.chapter4.triangulated-basics.g015
	\[\begin{tikzcd}
		X \arrow[r, "f"] & Y \\
		X' \arrow[r, "{f'}"'] \arrow[u, "u"] & Y' \arrow[u, "v"']
	\end{tikzcd}\]
	can be extended to a diagram
% segment-id: o014.li2.chapter4.triangulated-basics.g016
	\[\begin{tikzcd}
		TX' \arrow[r, "{Tf'}"] & TY' \arrow[r, "{Tg'}"] & TZ' \arrow[r, "{- Th'}"] & T^2 X' \arrow[phantom, ld, "\star" description] \\
		X'' \arrow[u, "o"] \arrow[r, "{f''}"] & Y'' \arrow[u, "p"] \arrow[r, "{g''}"] & Z'' \arrow[u, "q"] \arrow[r, "h''"] & TX'' \arrow[u, "{-To}"'] \\
		X \arrow[r, "f"] \arrow[u, "r"] & Y \arrow[r, "g"] \arrow[u, "s"] & Z \arrow[r, "h"] \arrow[u, "t"] & TX \arrow[u, "Tr"'] \\
		X' \arrow[r, "{f'}"'] \arrow[u, "u"] & Y' \arrow[r, "{g'}"'] \arrow[u, "v"] & Z' \arrow[r, "{h'}"'] \arrow[u, "w"] & TX' \arrow[u, "Tu"']
	\end{tikzcd}\]
	such that every row and column is a distinguished triangle, the square
	marked $\star$ is anticommutative (that is,
	$(To) h'' = -(Th') q$), and every other square is commutative.
\end{lemma}

% segment-id: o014.li2.chapter4.triangulated-basics.pf012
\begin{proof}
	We construct the required distinguished triangles step by step and then
	prove their commutativity and anticommutativity properties. First apply
	(TR2) separately to $u$, $v$, $f$, and $f'$ to obtain the portion
% segment-id: o014.li2.chapter4.triangulated-basics.g017
	\[\begin{tikzcd}[row sep=small, column sep=small]
		\bullet & \bullet & & \\
		\bullet \arrow[u] & \bullet \arrow[u] & & \\
		\bullet \arrow[r] \arrow[u] & \bullet \arrow[r] \arrow[u] & \bullet \arrow[r] & \bullet \\
		\bullet \arrow[r] \arrow[u] & \bullet \arrow[u] \arrow[r] & \bullet \arrow[r] & \bullet
	\end{tikzcd}\]
	of the diagram, with every row and column distinguished. Next apply (TR2)
	to $fu = vf' : X' \to Y$ to obtain a distinguished triangle
	$X' \to Y \xrightarrow{m} A \xrightarrow{n} TX'$. The following
	construction is based on (TR5).

% segment-id: o014.li2.chapter4.triangulated-basics.l003
	\begin{enumerate}
% segment-id: o014.li2.chapter4.triangulated-basics.i006
		\item Apply (TR5) to the three distinguished triangles containing
		$f'$, $v$, and $vf'$ to obtain the following morphism of
		distinguished triangles:
% segment-id: o014.li2.chapter4.triangulated-basics.g018
		\begin{equation}\label{eqn:triangulated-9-diag-aux-0}\begin{tikzcd}
			X' \arrow[r, "{f'}"] \arrow[equal, d] & Y' \arrow[r, "{g'}"] \arrow[d, "v"] & Z' \arrow[r, "{h'}"] \arrow[d, "i"] & TX' \arrow[equal, d] \\
			X' \arrow[r, "{vf'}"] \arrow[d, "{f'}"'] & Y \arrow[r, "m"] \arrow[equal, d] & A \arrow[r, "n"] \arrow[d, "j"] & TX' \arrow[d, "{Tf'}"] \\
			Y' \arrow[r, "v"] \arrow[d, "{g'}"'] & Y \arrow[r, "s"] \arrow[d, "m"] & Y'' \arrow[r, "p"] \arrow[equal, d] & TY' \arrow[d, "{Tg'}"] \\
			Z' \arrow[r, "i"'] & A \arrow[r, "j"'] & Y'' \arrow[r, "k"'] & TZ'
		\end{tikzcd}\end{equation}
		The last row is a newly constructed distinguished triangle.

% segment-id: o014.li2.chapter4.triangulated-basics.i007
		\item Apply (TR5) to the three distinguished triangles containing
		$u$, $f$, and $fu$ to obtain the following morphism of distinguished
		triangles:
% segment-id: o014.li2.chapter4.triangulated-basics.g019
		\begin{equation}\label{eqn:triangulated-9-diag-aux-1}\begin{tikzcd}
			X' \arrow[r, "u"] \arrow[equal, d] & X \arrow[r, "r"] \arrow[d, "f"] & X'' \arrow[r, "o"] \arrow[d, "a"] & TX' \arrow[equal, d] \\
			X' \arrow[r, "fu"] \arrow[d, "u"'] & Y \arrow[r, "m"] \arrow[equal, d] & A \arrow[r, "n"] \arrow[d, "b"] & TX' \arrow[d, "Tu"] \\
			X \arrow[r, "f"] \arrow[d, "r"'] & Y \arrow[r, "g"] \arrow[d, "m"] & Z \arrow[r, "h"] \arrow[equal, d] & TX \arrow[d, "Tr"] \\
			X'' \arrow[r, "a"'] & A \arrow[r, "b"'] & Z \arrow[r, "c"'] & TX''
		\end{tikzcd}\end{equation}
		The last row is a newly constructed distinguished triangle.

% segment-id: o014.li2.chapter4.triangulated-basics.i008
		\item Apply (TR2) to $f'' := ja: X'' \to Y''$ to obtain a
		distinguished triangle
		$X'' \xrightarrow{f''} Y'' \xrightarrow{g''} Z''
		\xrightarrow{h''} TX''$.

% segment-id: o014.li2.chapter4.triangulated-basics.i009
		\item Consider the following distinguished triangles obtained from
		the preceding steps and rotations (TR3):
% segment-id: o014.li2.chapter4.triangulated-basics.q008
		\begin{gather*}
			X'' \xrightarrow{a} A \xrightarrow{b} Z \xrightarrow{c} TX'', \quad A \xrightarrow{j} Y'' \xrightarrow{k} TZ' \xrightarrow{-Ti} TA, \\
			X'' \xrightarrow{f'' = ja} Y'' \xrightarrow{g''} Z'' \xrightarrow{h''} TX'' .
		\end{gather*}
		Applying (TR5) then gives a morphism of distinguished triangles
% segment-id: o014.li2.chapter4.triangulated-basics.g020
		\begin{equation}\label{eqn:triangulated-9-diag-aux-2}\begin{tikzcd}
			X'' \arrow[r, "a"] \arrow[equal, d] & A \arrow[r, "b"] \arrow[d, "j"] & Z \arrow[r, "c"] \arrow[d, "t"] & TX'' \arrow[equal, d] \\
			X'' \arrow[r, "{f''}"] \arrow[d, "a"'] & Y'' \arrow[r, "{g''}"] \arrow[equal, d] & Z'' \arrow[r, "{h''}"] \arrow[d, "q"] & TX'' \arrow[d, "Ta"] \\
			A \arrow[r, "j"] \arrow[d, "b"'] & Y'' \arrow[r, "k"] \arrow[d, "m"] & TZ' \arrow[r, "{-Ti}"] \arrow[equal, d] & TA \arrow[d, "Tb"] \\
			Z \arrow[r, "t"'] & Z'' \arrow[r, "q"'] & TZ' \arrow[r, "l"'] & TZ
		\end{tikzcd}\end{equation}
		The last row is a newly constructed distinguished triangle. Set
		$w := bi$. Commutativity of
		\eqref{eqn:triangulated-9-diag-aux-2} implies $l = -Tw$; rotation
		therefore gives the distinguished triangle
% segment-id: o014.li2.chapter4.triangulated-basics.q009
		\[ Z' \xrightarrow{w} Z \xrightarrow{t} Z'' \xrightarrow{q} TZ' . \]
	\end{enumerate}

% segment-id: o014.li2.chapter4.triangulated-basics.p009
	The operations above give the second row and third column of the required
	diagram. The remaining first row (or fourth column) is obtained by
	rotating the fourth row (or first column) according to (TR3), and is
	therefore also a distinguished triangle (see
	Remark~\ref{rem:triangle-signs}).

	We now verify the commutativity and anticommutativity properties. From
	$w = bi$ and $s = jm$, as implied by
	\eqref{eqn:triangulated-9-diag-aux-0}, it follows that $(u, v, w)$ and
	$(r, s, t)$ are respectively the composites
% segment-id: o014.li2.chapter4.triangulated-basics.q010
	\begin{gather*}
		(X', Y', Z') \xrightarrow[\eqref{eqn:triangulated-9-diag-aux-0}]{(\identity, v, i)} (X', Y, A) \xrightarrow[\eqref{eqn:triangulated-9-diag-aux-1}]{(u, \identity, b)} (X, Y, Z), \\
		(X, Y, Z) \xrightarrow[\eqref{eqn:triangulated-9-diag-aux-1}]{(r, m, \identity)} (X'', A, Z) \xrightarrow[\eqref{eqn:triangulated-9-diag-aux-2}]{(\identity, j, t)} (X'', Y'', Z''),
	\end{gather*}
	so both are morphisms of triangles. Finally, consider the morphism
	$(o, p, q, -To)$ in the diagram. It remains to prove that the following
	diagram is commutative:
% segment-id: o014.li2.chapter4.triangulated-basics.g021
	\[\begin{tikzcd}
		TX' \arrow[r, "{Tf'}"] & TY' \arrow[r, "{Tg'}"] & TZ' \arrow[r, "{- Th'}"] & T^2 X' \\
		A \arrow[u, "n"] \arrow[r, "j"] & Y'' \arrow[r, "k"] \arrow[u, "p"] & TZ' \arrow[equal, u] \arrow[r, "-Ti"] & TA \arrow[u, "Tn"'] \\
		X'' \arrow[u, "a"] \arrow[r, "{f''}"'] \arrow[bend left=60, uu, "o"] & Y'' \arrow[r, "{g''}"'] \arrow[equal, u] & Z'' \arrow[r, "{h''}"'] \arrow[u, "q"] & TX'' \arrow[u, "Ta"'] \arrow[bend right=60, uu, "To"'] .
	\end{tikzcd}\]
	Indeed, \eqref{eqn:triangulated-9-diag-aux-1} implies $o = na$, so both
	wings commute. Since the three squares on the right of
	\eqref{eqn:triangulated-9-diag-aux-0} commute, the upper part also
	commutes; commutativity of the lower part follows from
	\eqref{eqn:triangulated-9-diag-aux-2}. This proves the assertion.
\end{proof}

% segment-id: o014.li2.chapter4.triangulated-basics.p011
Except for Proposition~\ref{prop:adjoint-triangulated}, every assertion in
this section also holds for $\Bbbk$-linear categories, where $\Bbbk$ is a
commutative ring.

% segment-id: o014.li2.chapter4.triangulated-basics.r002
\begin{remark}
	\index{compact object!in a triangulated category}
	\index{generator!in a triangulated category}
	\index{Brown Representability Theorem}
	Generators and compact objects in triangulated categories are important
	topics not treated in this book. Although the ideas resemble the versions
	in Definition~\ref{def:generators} and
	\sourcecrossref{def:cpt-objects}{A.2.2}, these notions require different
	definitions in a triangulated category. This leads to an important form
	of the Brown Representability Theorem for triangulated categories: if
	$\mathcal{D}$ is a “compactly generated” triangulated category with small
	$\coprod$, and a cohomological functor
	$H: \mathcal{D}^{\opp} \to \cate{Ab}$ maps small $\coprod$ in
	$\mathcal{D}$ to $\prod$ in $\cate{Ab}$, then $H$ is representable.
	Readers who need this result are advised to consult
	\cite[Tags 09SI, 09SM, 0A8E]{stacks}.

	Notice the family resemblance between the Brown Representability Theorem
	and the various representability and adjoint-functor theorems in
	\S\ref{sec:AFT}, \S\ref{sec:Grothendieck-cat}, and
	\S~\sourcecrossref{sec:cpt-objects}{A.2}.
\end{remark}
